\section{Reversal balance and a finite-range inverse}
\label{sec:reversal}
A spatial first-impact law may reveal the boundary directly through its
support. We ask instead what remains when the output of each attempt is
only a collision bit. The starting distribution is now a controlled
input. The distinction matters: reducing the recorded output while
allowing localized preparations does not establish dominance over a
uniform-launch experiment. Our principal identity separates the two
issues. It expresses an occupation difference by two scalar collision
probabilities at the same positive horizon, with no small-time limit.

Let $\mathcal O\subset\R^d$ be closed, and put $\chi=\one_{\mathcal O}$.
An attempt starts at $q$ with velocity $v$, $|v|=1$, and runs until time
$t>0$. A start in $\mathcal O$ is an unsuccessful preparation and is
assigned output zero, as is a free start with no collision by time $t$.
A free start with a collision by that time has output one. The two zero
outcomes are not distinguished by the observation. The event includes
the endpoint time; boundary conventions are immaterial for the smooth
classes and absolutely continuous preparations below. Write $a=tv$ and
\begin{equation}\label{eq:bit}
 B_a(q)=\one_{\{q\notin\mathcal O,\ [q,q+a]\cap\mathcal O\ne\varnothing\}}.
\end{equation}
The straight segment in this definition is the unreflected path. It
detects the physical first collision, because the trajectory is straight
until that collision; no part of the path after the first collision is
observed. For a nonnegative probability density $f$ set
\[
 P_a(f)=\int_{\R^d}f(q)B_a(q)\dd q,\qquad f_a(q)=f(q-a).
\]
The denominator is the number of attempted preparations, not the number
of free or successful starts. The experiment commanding $f_a$ shifts
the starting distribution by $+a$ and reverses the velocity. It does
not run a time-reversed reflected trajectory.

\begin{theorem}[Reversal balance]\label{thm:balance}
For every closed obstacle set and every displacement $a$,
\begin{align}
 B_a(q)-B_{-a}(q+a)&=\chi(q+a)-\chi(q),\label{eq:balance-point}\\
 P_a(f)-P_{-a}(f_a)&=
       \int f(q)\{\chi(q+a)-\chi(q)\}\dd q.\label{eq:balance}
\end{align}
The pointwise identity requires no convexity, separation or periodicity.
It remains valid when the segment intersects several components.
\end{theorem}
\begin{proof}
Let $I(q)$ indicate whether $[q,q+a]$ meets $\mathcal O$. The two
oriented segments have the same image, so the left side of
\eqref{eq:balance-point} is
$\{\chi(q+a)-\chi(q)\}I(q)$. If either endpoint is in $\mathcal O$,
then $I(q)=1$; if neither is, the prefactor is zero. This proves the
pointwise assertion. In the reverse probability change variables from
$q+a$ to $q$. The result is \eqref{eq:balance}.
\end{proof}

In particular, the pair of functionals $P_a,P_{-a}$ on smooth compactly
supported launch densities determines the signed measure with density
$\delta_a\chi(q):=\chi(q+a)-\chi(q)$. Testing all nonnegative smooth
densities determines this measure uniquely, by scaling and linearity.
These are functionals of a controllable spatial input, not two real
observations. For finite acquisition we use an explicit finite family
of inputs rather than appeal to exact differentiation of such a measure.

A difference operator ordinarily leaves an additive constant along each
translation orbit undetermined. The next theorem fixes it by a geometric
zero condition, without providing a free reference point.

\begin{theorem}[Finite-chain inversion]\label{thm:chain}
Let $u\ge0$ be a function, $a\ne0$, and $m\ge1$. Suppose that for every
$x$ in a region $U$,
\begin{equation}\label{eq:zero-chain}
                 \min_{0\le k\le m}u(x+ka)=0.
\end{equation}
Then the values of $\delta_a u$ on
$U^+=\bigcup_{j=0}^{m-1}(U+ja)$ determine $u$ on $U$ by
\begin{equation}\label{eq:min-inverse}
 u(x)=-\min_{0\le k\le m}\sum_{j=0}^{k-1}\delta_a u(x+ja),
\end{equation}
where the empty sum is zero. If $u'$ obeys the same zero condition,
then, for $1\le p\le\infty$,
\begin{equation}\label{eq:chain-stability}
 \|u-u'\|_{L^p(U)}\le
       m\|\delta_a u-\delta_a u'\|_{L^p(U^+)}.
\end{equation}
More generally, perturbing every difference in a chain by at most $b$
perturbs the reconstruction by at most $mb$.
\end{theorem}
\begin{proof}
The partial sum is $u(x+ka)-u(x)$. Its minimum equals $-u(x)$ by
\eqref{eq:zero-chain}, proving the formula. The minimum of two finite
lists differs by at most their maximum entrywise difference. Hence
\[
 |u(x)-u'(x)|\le
   \sum_{j=0}^{m-1}|\delta_a u(x+ja)-\delta_a u'(x+ja)|.
\]
Minkowski's inequality and translation of each integration domain prove
\eqref{eq:chain-stability}; the uniform perturbation statement follows
from the same inequality. Almost-everywhere hypotheses suffice: only
finitely many translates of exceptional null sets occur.
\end{proof}

\begin{lemma}[A geometric zero witness]\label{lem:zero}
Suppose $\mathcal O$ is a union of compact components of diameter at
most $D_0$, with mutual distance at least $d_0>0$. If
$0<t<d_0$, $a=tv$ and $mt>D_0$, then $u=\chi$ satisfies
\eqref{eq:zero-chain} everywhere.
Fix $\sigma_0>0$ with $t<d_0-2\sigma_0$ and
$mt>D_0+2\sigma_0$. If $\kappa\ge0$ has integral one and support
in the closed unit ball, then for every $0<\sigma\le\sigma_0$ the
local occupation average
\begin{equation}\label{eq:occupation}
 u_\sigma(x)=\int \sigma^{-d}\kappa((q-x)/\sigma)\chi(q)\dd q
\end{equation}
also satisfies \eqref{eq:zero-chain}.
\end{lemma}
\begin{proof}
If all chain points belong to $\mathcal O$, successive points cannot
belong to different components, since their distance is $t<d_0$.
They therefore lie in one component, contradicting the endpoint distance
$mt>D_0$. For the second assertion, $u_\sigma(x)>0$ implies
$x\in\mathcal O+\sigma\overline B$. These enlarged components have
diameter at most $D_0+2\sigma$ and separation at least $d_0-2\sigma$.
The same argument excludes positivity at every point of the chain.
\end{proof}

\begin{corollary}[Local rigidity from scalar launch laws]\label{cor:local}
Under the diameter and separation hypotheses, the two probability
functionals at one fixed horizon and opposite directions determine the
obstacle indicator almost everywhere on $U$, using only launch densities
supported on its bounded chain enlargement. For regular closed bodies
they determine the bodies themselves in the interior of $U$.
Alternatively, the same determination follows from localized kernel
preparations at all sufficiently small scales, with the explicit
finite-chain formula at each scale.
\end{corollary}
\begin{proof}
Apply Theorem~\ref{thm:balance} to recover $\delta_a\chi$ and
Theorem~\ref{thm:chain} with Lemma~\ref{lem:zero}. Two different regular
closed convex unions would differ on an open set, so almost-everywhere
equality determines the set. For the alternative, the balance identity
recovers $\delta_a u_\sigma$ exactly, the minimum formula recovers
$u_\sigma$, and approximate-identity convergence gives $\chi$ almost
everywhere as $\sigma\downarrow0$.
\end{proof}

The zero witness is substantive. For constant sequences both zero and
one have zero difference, and for unbounded slabs parallel to $a$ the
same ambiguity can persist. Here the witness is implied by bounds on
physical component size and separation, not by an unobserved boundary
value. The local theorem uses no periodicity or boundary smoothness.
In the convex billiard class and at $t<d_0$, there cannot be a second
collision: a reflected exterior ray cannot re-enter the same convex
body, and another body is at least $d_0$ away. Thus the successful
output is also the one-collision count at this horizon, with solid
preparations still counted as zero attempts.
